\section{Conclusion}
\label{sec:conclusion}

This work presented and evaluated CloudAware-MOCCA, a context-aware
self-adaptive middleware for mobile cloud and edge offloading, and its results
yield three findings that go beyond the specific prototype. First, the physical
deployment demonstrates that integrated adaptive offloading is feasible in
practice rather than only in simulation. A complete monitor-analyze-plan-execute
control loop, running on a commodity Android device against containerized edge
and cloud tiers with transparent server-side forwarding under load, sustained
real task execution with a low fallback rate and consistently avoided the
penalty of a fixed always-cloud policy, confirming that the architecture behaves
as intended on real hardware. Second, a compact learned model can reproduce a
hand-tuned deterministic policy with high fidelity, but such fidelity must not be
mistaken for independent decision superiority. Because the model is trained on
the rule engine's own decisions, it faithfully imitates them at negligible
on-device inference cost, yet on a label-independent quality metric it performs
only on par with the policy it copies; demonstrating genuine improvement would
require learning from measured execution outcomes rather than from
rule-generated labels. Third, energy-model calibration and remote execution
duration emerge as the principal limitations for energy-aware offloading, since
pricing only radio activity rather than whole-device power across the full round
trip caused a substantial share of remote executions to cost more energy than
local execution. Addressing these limitations through duration-aware and
payload-aware cost models, per-device power calibration, and quantified runtime
overhead constitutes the most promising direction for future work.
